\documentclass[../../../include/open-logic-section]{subfiles}

\begin{document}

\olfileid[id]{sth}{choice}{justifications}
\olsection{Pertimbangan Intrinsik tentang Pilihan}

Pertanyaan yang lebih luas ialah apakah Pengurutan Baik, Pilihan, atau bahkan
keterbandingan semua himpunan menurut ukurannya---tidak menjadi soal yang
mana---dapat dibenarkan.

Berikut suatu upaya pembenaran \emph{intrinsik}. Sebelumnya, dalam
\olref[z][story]{sec}, kita memperkenalkan beberapa prinsip tentang
hierarki. Salah satunya patut dinyatakan kembali:
\begin{enumerate}
	\item[] \stagesacc. Untuk sebarang tahap $S$ dan sebarang himpunan yang
	terbentuk \emph{sebelum} tahap $S$: pada tahap $S$ terbentuk suatu
	himpunan yang anggota-anggotanya tepat merupakan himpunan-himpunan
	tersebut. Tidak ada hal lain yang terbentuk pada tahap~$S$.
\end{enumerate}
Sesungguhnya, banyak penulis telah mengemukakan bahwa Aksioma Pilihan dapat
dibenarkan melalui prinsip ini (atau sesuatu yang serupa dengannya). Kita
akan memberikan suatu uraian ringkas mengenai pendekatan tersebut.

Kita mulai dengan suatu hasil sederhana yang menawarkan \emph{satu lagi}
perumusan ekuivalen bagi Pilihan:

\begin{thm}[dalam $\ZF$]\ollabel{choiceset}
Pilihan ekuivalen dengan prinsip berikut. Jika !!{element}s dari $A$ saling
lepas dan tak kosong, maka terdapat suatu $C$ sedemikian sehingga
$C\cap x$ merupakan singleton bagi setiap $x\in A$. (Kita menyebut $C$
seperti ini sebagai suatu himpunan pilihan bagi~$A$.)
\end{thm}

Bukti hasil ini lugas dan kita serahkan sebagai latihan kepada pembaca.

\begin{prob}
Buktikan \olref[sth][choice][justifications]{choiceset}. Jika Anda mengalami
kesulitan, suatu bukti dapat ditemukan dalam
\cite[hlm.~242--243]{Potter2004}.
\end{prob}

Pokok yang mendasar ialah bahwa himpunan pilihan bagi $A$ tidak lain adalah
daerah hasil suatu fungsi pilihan bagi~$A$. Jadi, untuk membenarkan Pilihan,
kita cukup mencoba membenarkan perumusan ekuivalennya dalam hal keberadaan
himpunan pilihan. Itulah yang sekarang akan kita coba lakukan.

Andaikan !!{element}s dari $A$ saling lepas dan tak kosong. Berdasarkan
\stageshier{} (lihat \olref[z][story]{sec}), $A$ terbentuk pada suatu
tahap~$S$. Perhatikan bahwa semua !!{element}s dari $\bigcup A$ tersedia
sebelum tahap~$S$. Berdasarkan \stagesacc{}, bagi \emph{sebarang} himpunan
yang terbentuk sebelum~$S$, terbentuk pula suatu himpunan yang
anggota-anggotanya tepat merupakan himpunan-himpunan tersebut. Dengan kata
lain, setiap koleksi \emph{yang mungkin} dari himpunan-himpunan yang tersedia
sebelumnya akan ada pada tahap~$S$. Namun, tentu \emph{mungkin} untuk memilih
objek-objek yang dapat membentuk suatu himpunan pilihan bagi~$A$; ini hanya
merupakan suatu himpunan bagian yang sangat spesifik dari $\bigcup A$.
Dengan demikian, himpunan pilihan seperti itu ada, sebagaimana diperlukan.

Itulah suatu upaya yang \emph{sangat} singkat untuk membenarkan Pilihan
berdasarkan pertimbangan intrinsik. Namun, untuk mendalami gagasan ini, Anda
sebaiknya membaca pengembangan yang rapi oleh Potter
(\citeyear[\S14.8]{Potter2004}).

\end{document}
